\section{Unsigned cancellation and its geometric input}
\label{sec:model}

\subsection{The finite decision problem}

Let $R$ be a finite set of row labels and $C$ a finite set of columns. The
matrix $A=(a_{rc})$ has nonnegative integer entries and is stored sparsely.
Each row either labels one edge of a finite connected undirected
multigraph $G$ or is declared unavailable for top cancellation. The graph
has a distinguished exterior vertex $0$. Parallel edges are distinct row
labels; loops cannot belong to a forest. Write $E\subseteq R$ for the
eligible nonloop edges and $r=|V(G)|-1$ for the rank.

A top choice is a spanning tree $T\subseteq E$. It removes the rows of
$T$. On the remaining row set $L$, a column $c$ is a legal leaf when
\[
                         \sum_{r\in L}a_{rc}=1.
\]
There is then a unique row $u\in L$ with $a_{uc}=1$; delete $u$ and
record the pair $(u,c)$. Continue until every row is deleted or no leaf
remains. The input question is whether \emph{some} top tree gives complete
peeling. Rows unavailable for top cancellation can still be deleted by
leaf pivots; excluding them from the matrix would change the problem.

\begin{lemma}[Saturation preserves the finite problem]
\label{lem:saturation}
Replacing every positive entry $a$ by $\min(2,a)$ preserves the legal
leaf predicate on every subset of rows, and therefore preserves every
peeling trace and the answer for every top tree.
\end{lemma}
\begin{proof}
A sum of nonnegative integers is one exactly when one summand is one and
every other summand is zero. Saturation preserves zero, one, and the
property of being at least two. This applies to each column on every
row subset.
\end{proof}

This observation does \emph{not} license saturating counts before a
parity or signed-chain calculation: an odd count of three would become
even. The implementation's input digest binds the uncapped integers,
and its cycle-compatibility check uses their actual parities.

\begin{definition}[Stopping set]
A nonempty subset $S\subseteq R$ is a stopping set if every column sum
$\sum_{r\in S}a_{rc}$ is zero or at least two.
\end{definition}

\begin{proposition}[Exact immutable obstruction]
\label{prop:stopping}
Peeling a row subset $L$ completes if and only if $L$ contains no stopping
set. Its final residue is independent of the order of legal leaf choices
and equals the union of all stopping subsets of $L$.
\end{proposition}
\begin{proof}
A stopping set cannot lose its first row to a legal leaf: the corresponding
column would already have sum at least two inside that stopping set.
Consequently every stopping subset remains throughout every peeling
execution. Conversely, a nonempty terminal residue has no unit column and
is itself stopping. The union of stopping sets is stopping: for a column
with a positive contribution to the union, one constituent stopping set
already contributes at least two. The union of all stopping subsets is therefore the largest stopping
residue, possibly empty. Every terminal residue contains this union, and a
nonempty terminal residue is itself stopping and hence contained in it.
If the union is empty, all executions complete.
\end{proof}

The preceding incoming reports prove that complete positive raw
cancellation can always be witnessed by such immutable column peeling
\cite{repo-positive,repo-dual-normalization}. This completion theorem is
the reason a search over more complicated positive substitutions is
unnecessary for the present decision problem. A certificate below uses
only the original matrix entries and an ordinary tree.

\subsection{Activation reverses the viewpoint}

Call a removed row \emph{active}. A top tree is the initial active set.
A column with support $H$ can activate its row $u$ exactly when
\[
               a_{uc}=1,\qquad H\setminus\{u\}\subseteq X,
\]
where $X$ is the current active set. Thus it gives the Horn rule
$H\setminus\{u\}\to u$ for each of its unit entries. A column with no
unit entry gives no rule. A unit singleton is an unconditional rule;
two-row supports give directed unary rules. Supports of three or more
rows are the nonlinear part of the system.

The closure $\cl_A(X)$ is obtained by repeatedly applying all legal
rules. Finite monotonicity shows that it exists and is independent of the
application order. The tree $T$ is successful exactly when
$\cl_A(T)=R$. This is an independent-key problem for a Horn closure,
where independence is taken in an external graphic matroid.

\subsection{Source-relative geometric interpretation}

The native complementary core is a handle structure obtained from an
authenticated triangulation and normal coordinate vector. After fixing
the lower $0/1$ forest, its surviving one-handle cocores index columns.
Two-handle attaching curves give \emph{geometric occurrence counts};
algebraic exponent sums cannot replace them. The cap graph records the
two actual three-handle ends of each source two-handle, with exposed ends
sent to the exterior root. Its rooted spanning trees give legal $2/3$
cancellation schedules under the source-relative theorem proved in the
incoming packages \cite{repo-dual-normalization,repo-relative-top}.

After top cancellation, an immutable column leaf supplies a one/two
cancellation in the already authenticated source model. Finite bundle
boundary relations must also be included. In particular, later annular
rows are not eligible cap edges and are represented as unavailable rows.
Their omission would turn an incomplete core criterion into an unsound
whole-cut criterion.

The graph--matrix algorithm does not independently certify that these
data arise from a manifold. Its positive output becomes a geometric
certificate only after the cap graph, the occurrence matrix, and the
required relative charts have been reconstructed and checked. A negative
output says that this \emph{fixed raw forest-and-peeling criterion} cannot
complete. It does not say that the cut piece is not a handlebody, or that
the original knot is nontrivial.

\subsection{Why the exceptional column shape is natural}

For original native one-handle cocores, a triangular source-face region
has at most three actual boundary occurrences. A column containing a
unit and more than two owners therefore has precisely three distinct
unit owners. Moreover, the raw parity incidences obey the
boundary-of-boundary identity \cite{repo-dual-normalization}. The new
source audit rechecks these properties before saturation. It also finds
the same three-occurrence bound on every reconstructed finite boundary
matrix in the maintained corpus. The latter is an experimental statement
about those matrices, not a universal theorem about arbitrary interfaces.

The parameter $k$ will count nonlinear columns after free cancellations.
The local theorem supplies no uniform polylogarithmic bound on $k$ in
the original crossing number, and does not derive one from the underlying
manifold type or from the number of bits per normal coordinate. Establishing useful
geometric control of $k$ is one of the main next questions.
